\section{Conclusion}
\label{sec:conclusion}

We have introduced ApportionRegions (AR), an algorithm that
generates congressional district maps by applying the prime factorization of
the apportioned seat count as a hierarchical geographic partition. The
algorithm has two inputs: arithmetic (the prime factorization of $n$) and
geography (the census-tract adjacency graph with TIGER boundary weights).
Under a fixed implementation, seed rule, and tie-breaking policy, it produces
a reproducible map for each state at each seat count, with no political inputs
or manual boundary choices.

The key structural properties --- reuse of canonical prime splits across
seat counts, stability under reapportionment when factorizations share a
common prefix, and complete structural breaks when they do not --- follow
directly from the Fundamental Theorem of Arithmetic. The tree topology is
therefore arithmetic rather than discretionary, while the exact geography still
depends on the published graph, seed, and METIS realization.

The constitutional argument is that AR can be read as extending
Huntington-Hill (Article~I~\S2) from seat allocation to geographic allocation,
rather than treating the Manner clause (\S4) as the only possible authority.
The apportionment process, taken to its geographic conclusion, supplies the
tree structure that the partitioner realizes.

The most striking empirical finding is the $51 \to 52$ California transition,
where consecutive seat counts have completely disjoint prime structures,
forcing a full map redesign. This is not a defect of AR --- it is a
consequence of the mathematics. Any redistricting algorithm that takes the
arithmetic of seat counts seriously will face the same transitions. AR makes
them visible and principled.

Future work: seed sensitivity of large-prime fallback splits;
formal proof of uniqueness for the canonical split prescription; empirical
disruption study for 2010$\to$2020 reapportionment transitions;
and the AR-smooth variant for minimising tract reassignments.
